\motto{I speak twelve languages: English is the \emph{bestest}.

--- Stefan Bergman\footnote{Stefan Bergman (1895--1977), bearing a very Scandinavian name, was a Polish-American mathematician best known for his work in complex analysis, especially in several complex variables. He introduced the Bergman kernel, a fundamental concept in complex analysis that has influenced many areas of mathematics and theoretical physics.

Bergman was born in Poland (then part of the Russian Empire), and studied in Berlin. He fled Europe during World War II and eventually settled in the United States.}}
\chapter{Partial Bergman kernels}
\label{chap:Bergman}

The classical Bergman kernel of a Hermitian line bundle $L$ on a compact Kähler manifold $X$ encodes, through its asymptotic expansion as the tensor power $L^k$ goes to infinity, deep geometric information about the underlying metric. This chapter develops the corresponding \emph{partial} Bergman kernel. The word partial refers to the fact that we only use sections cut out by an ideal, typically the multiplier ideal of an auxiliary psh weight on a closed subset $K\subseteq X$.

The main result is the convergence theorem \cref{thm:pBMconvergence}. Given a $\theta$-psh potential $\varphi\in \PSH(X,\theta)_{>0}$, a closed non-pluripolar weighted subset $(K,v)$, and a Bernstein--Markov measure $\mu$, the normalized partial Bergman measures on $X$ converge weakly to the partial equilibrium measure of $(K,v,\varphi)$ as $k\to\infty$. This is the relative version of the theorem of Berman--Boucksom--Witt Nystr\"om \cite{BBWN11}, corresponding to the case $\varphi=V_\theta$. The assumptions are essentially the mildest ones under which the statement makes sense.

The first part of the chapter introduces the relative envelope operators. In \cref{sec:envrel} we define the relative $P$-envelope $P_{\theta,K}[\varphi](v)$ and the relative $\mathcal{I}$-envelope $P_{\theta,K}[\varphi]_{\mathcal{I}}(v)$ by replacing the constraint $\eta\leq 0$ in the usual $P$- and $\mathcal{I}$-envelopes with the constraint $\eta|_K\leq v$. We establish their basic properties, including monotonicity, mass-preservation, and the relation with the envelope operators of \cref{chap:envelopes}.

The second part develops the equilibrium theory associated with a weighted subset $(K,v)$. In \cref{sec:pem} we show that the Monge--Amp\`ere measure $\theta_{P_{\theta,K}[\varphi](v)}^n$ is carried by the contact locus $\{P_{\theta,K}[\varphi](v)|_K=v\}\subseteq K$, which justifies the name \emph{partial equilibrium measure}. The chapter concludes in \cref{sec:quant} with the proof of the weak convergence of partial Bergman measures towards this partial equilibrium.

\section{Partial envelopes}\label{sec:envrel}
Let $X$ be a connected compact Kähler manifold of dimension $n$ and $K \subseteq X$ be a closed non-pluripolar set.
Let $\theta$ be a smooth closed real $(1,1)$-form on $X$ representing a pseudo-effective cohomology class.

\begin{definition}\label{def:partial_env}
    Given a function $f\colon K\rightarrow [-\infty,\infty]$ and $\varphi\in \PSH(X,\theta)$, we introduce the \emph{relative $P$-envelope}\index{envelope!P-envelope@$P$-envelope!relative P-envelope@relative $P$-envelope} of $\varphi$ (with respect to $K,f,\theta$) as
    \begin{equation}\label{eq:PthetaKvarphiv}
    P_{\theta,K}[\varphi](f)\coloneqq  \uscsup \Bigl\{ \eta \in \PSH(X,\theta): \eta|_K \leq f  \textup{ quasi-everywhere and } \eta\preceq \varphi\Bigr\}.
    \end{equation}
    \index{$P_{\theta,K}[\varphi](f)$}

    Similarly, we define the \emph{relative $\mathcal{I}$-envelope}\index{envelope!I-envelope@$\mathcal{I}$-envelope!relative I-envelope@relative $\mathcal{I}$-envelope} of $\varphi$ (with respect to $K,f,\theta$) as
    \begin{equation}\label{eq:PthetaKvarphivI}
    P_{\theta,K}[\varphi]_{\mathcal{I}}(f)\coloneqq  \uscsup \Bigl\{ \eta \in \PSH(X,\theta): \eta|_K \leq f  \textup{ quasi-everywhere and } \eta\preceq_{\mathcal{I}} \varphi\Bigr\}.
    \end{equation}
\end{definition}
By $\eta|_K\leq f$ quasi-everywhere, we mean that there is a pluripolar set $S\subseteq X$ so that $\eta|_{K\setminus S}\leq f|_{K\setminus S}$. In general the upper envelope in \eqref{eq:PthetaKvarphiv} or \eqref{eq:PthetaKvarphivI} may be identically $\infty$.

Observe that when $f$ is bounded from below, neither envelope is identically $-\infty$.



The envelopes contain essentially all envelope operators that we have studied in this book:
\begin{example}\label{ex:rel_env_cases}
    When $K=X$ and $f=0$, we have
    \[
        P_{\theta,X}[\varphi](0)=P_{\theta}[\varphi],\quad P_{\theta,X}[\varphi]_{\mathcal{I}}(0)=P_{\theta}[\varphi]_{\mathcal{I}}.
    \]
    When $K=X$, we have
    \[
        P_{\theta,X}[\varphi](f)=P_{\theta}[\varphi](f).
    \]
    These definitions reduce to the envelopes studied in \cref{chap:envelopes}.
    On the other hand, when $\varphi=V_{\theta}$, we have
    \[
    \begin{aligned}
        P_{\theta,K}[V_{\theta}](f)
        &=P_{\theta,K}[V_{\theta}]_{\mathcal{I}}(f)\\
        &=\uscsup \left\{ \eta \in \PSH(X,\theta):
        \begin{array}{l}
        \eta|_K \leq f \textup{ quasi-everywhere}
        \end{array}
        \right\}.
    \end{aligned}
    \]
    In an even more special case where $\varphi=V_{\theta}$ and $K=X$, we find that
    \[
        P_{\theta,X}[V_{\theta}](f)=\uscsup \Bigl\{ \eta \in \PSH(X,\theta): \eta \leq f \textup{ quasi-everywhere}\Bigr\}.
    \]
    This is nothing but $P_{\theta}(f)$ as we studied in \cref{def:Pthetaf}.
\end{example}



\begin{proposition}\label{prop:partial_env}
    Let $f,g\colon K\rightarrow [-\infty,\infty]$ be functions and $\varphi,\psi\in \PSH(X,\theta)$. Then we have the following:
    \begin{enumerate}
        \item\label{itm:prop:partial_env:1} We have
        \[
        P_{\theta,K}[\varphi](f)\leq P_{\theta,K}[\varphi]_{\mathcal{I}}(f);
        \]
        \item\label{itm:prop:partial_env:2} assume that $f$ is bounded from above, then
        \[
        \begin{aligned}
            \left.P_{\theta,K}[\varphi](f)\right|_K\leq & f\textup{ quasi-everywhere},\\
             \left. P_{\theta,K}[\varphi]_{\mathcal{I}}(f)\right|_K\leq & f\textup{ quasi-everywhere};
        \end{aligned}
        \]
        \item\label{itm:prop:partial_env:3} suppose that $\varphi\preceq_{\mathcal{I}}\psi$, then
        \[
            P_{\theta,K}[\varphi]_{\mathcal{I}}(f)\leq P_{\theta,K}[\psi]_{\mathcal{I}}(f);
        \]
        \item\label{itm:prop:partial_env:4} if $f\leq g$, then
        \[
        P_{\theta,K}[\varphi](f)\leq P_{\theta,K}[\varphi](g),\quad P_{\theta,K}[\varphi]_{\mathcal{I}}(f)\leq P_{\theta,K}[\varphi]_{\mathcal{I}}(g);
        \]
        \item\label{itm:prop:partial_env:5} for any $t \in (0,1)$, we have
        \[
        \begin{aligned}
            P_{\theta,K}\left[t\varphi+(1-t)\psi\right](f)\geq & tP_{\theta,K}[\varphi](f)+(1-t)P_{\theta,K}[\psi](f),\\
            P_{\theta,K}\left[t\varphi+(1-t)\psi\right]_{\mathcal{I}}(f)\geq & tP_{\theta,K}[\varphi]_{\mathcal{I}}(f)+(1-t)P_{\theta,K}[\psi]_{\mathcal{I}}(f).
        \end{aligned}
        \]
    \end{enumerate}
\end{proposition}
\begin{proof}
    \ref{itm:prop:partial_env:1} This follows from $\eta\preceq \varphi\Rightarrow \eta\preceq_{\mathcal I}\varphi$.

    \ref{itm:prop:partial_env:2} This follows from \cref{prop:supsupstardiff} and \cref{rmk:L1cpt}.

    \ref{itm:prop:partial_env:3} This follows by inclusion of the defining candidate sets.

    \ref{itm:prop:partial_env:4} This follows by inclusion of the defining candidate sets.

    \ref{itm:prop:partial_env:5} If $\eta,\eta'\in \PSH(X,\theta)$ are such that
    \[
    \begin{split}
        &\eta|_K\leq f\textup{ quasi-everywhere},\quad \eta'|_K\leq f\textup{ quasi-everywhere},\\
        &\eta\preceq \varphi,\quad \eta'\preceq \psi,
    \end{split}
    \]
    then for $t\in (0,1)$,
    \[
        \left.\left(t\eta+(1-t)\eta'\right)\right|_K\leq f\textup{ quasi-everywhere},\quad t\eta+(1-t)\eta' \preceq t\varphi+(1-t)\psi.
    \]
    Therefore,
    \[
    t\eta+(1-t)\eta'\leq P_{\theta,K}\left[t\varphi+(1-t)\psi\right](f).
    \]
    Since $\eta$ and $\eta'$ are arbitrary, the first inequality follows. The second follows from the same argument.
\end{proof}

Similar to the usual envelopes studied in \cref{chap:envelopes}, the $P$-partial envelopes are well-behaved only when $\varphi$ has positive mass.

\begin{lemma}\label{lma:pos_mass_partial}
    Suppose that $f\colon K\rightarrow (-\infty,\infty]$ is bounded from below, and $\varphi\in \PSH(X,\theta)_{>0}$. Then
    \begin{equation}\label{eq:Ptheta_partial_pos_mass}
    P_{\theta,K}[\varphi](f)= \uscsup \Bigl\{ \eta \in \PSH(X,\theta)_{>0}: \eta|_K \leq f  \textup{ quasi-everywhere and } \eta\preceq \varphi\Bigr\}.
    \end{equation}
\end{lemma}
In other words, in this situation, in the definition \eqref{eq:PthetaKvarphiv}, we may restrict to only those $\eta$ with positive mass.
\begin{proof}
    By our assumption, we can find $C>0$ so that $\varphi|_K-C\leq f$. Then if $\eta$ is a candidate in \eqref{eq:PthetaKvarphiv}, $\eta\lor (\varphi-C)$ is also a candidate, and $\eta\lor (\varphi-C)\in \PSH(X,\theta)_{>0}$. Therefore, $\eta\lor (\varphi-C)$ is also a candidate for the right-hand side of \eqref{eq:Ptheta_partial_pos_mass}. The equality \eqref{eq:Ptheta_partial_pos_mass} follows.
\end{proof}

Note that in the most important cases, we can actually replace the condition $\eta\leq f$ quasi-everywhere in \eqref{eq:PthetaKvarphiv} by everywhere.
\begin{lemma}\label{lma:PKoutsidepps}
    Let $f\colon K\rightarrow (-\infty,\infty]$ be bounded from below, and $\varphi\in \PSH(X,\theta)_{>0}$.
    Then
    \begin{equation}\label{eq:PthetaK_new_def}
        \begin{split}
    P_{\theta,K}[\varphi](f)=&\uscsup \Bigl\{ \eta \in \PSH(X,\theta): \eta|_K \leq f  \textup{ and } \eta\preceq \varphi\Bigr\}\\
    =& \uscsup \Bigl\{ \eta \in \PSH(X,\theta)_{>0}: \eta|_K \leq f  \textup{ and } \eta\preceq \varphi\Bigr\}.
        \end{split}
    \end{equation}
\end{lemma}
\begin{proof}
Both $\geq$ directions in \eqref{eq:PthetaK_new_def} are trivial.

Conversely, let $\eta$ be a candidate in \eqref{eq:Ptheta_partial_pos_mass}.
Let $S\subseteq X$ be a pluripolar set so that $\eta|_{K\setminus S}\leq f|_{K\setminus S}$.

We now claim that we can choose $\chi\in \PSH(X,\theta)_{>0}$ with $\chi|_S\equiv -\infty$ and $\chi\leq \eta$.

Assume this for the moment. For any $\delta\in (0,1)$, we have
\[
(1-\delta)\eta|_K+\delta \chi|_K\leq f,\quad (1-\delta)\eta+\delta \chi\preceq \varphi.
\]
Hence, $(1-\delta)\eta+\delta \chi$ is dominated by the right-hand side of \eqref{eq:PthetaK_new_def}.
Letting $\delta \to 0+$, we conclude that $\eta$ is dominated by the right-hand side of \eqref{eq:PthetaK_new_def} and hence, \eqref{eq:PthetaK_new_def} follows.

    It remains to construct $\chi$.
    Since $\varphi$ has positive mass, the class $\{\theta\}$ is big. Fix a K\"ahler form $\omega_0$ on $X$ and choose $\rho\in \PSH(X,\theta)$ with $\theta_{\rho}\geq 2\epsilon_0\omega_0$ for some $\epsilon_0>0$. By \cref{thm:gloJosefson}, there is a quasi-psh function $q$ such that $q|_S\equiv -\infty$. Choose $A>0$ so that $A\omega_0+\ddc q\geq 0$. For $a>0$ small enough, $\rho+aq\in \PSH(X,\theta)$ and still has pole $-\infty$ on $S$. Replacing $\rho$ by $\rho+aq$, we may therefore assume that $\rho|_S\equiv -\infty$.
    Since $\int_X \theta_{\eta}^n>0$, fixing such a $\rho$ and taking $\epsilon\in (0,1)$ small enough, we have
    \[
    \int_X \theta^n_{\epsilon \rho+(1-\epsilon)V_{\theta}}+\int_X \theta_{\eta}^n>\int_X \theta_{V_{\theta}}^n.
    \]
    Thus, by \cref{prop:landfinitecond1} and \cref{thm:diamond}, we have
    \[
    \Bigl(\epsilon \rho+(1-\epsilon)V_{\theta}\Bigr) \land \eta \in \PSH(X,\theta)_{>0}.
    \]
    It suffices to define this function as $\chi$.
\end{proof}
\begin{corollary}\label{cor:Ppartbdd_bdabove}
    Suppose that $f\colon K\rightarrow \mathbb{R}$ is bounded, and $\varphi\in \PSH(X,\theta)_{>0}$. Then there is a constant $C=C(\sup_K f,K,X)>0$ such that
        \[
            P_{\theta,K}[\varphi](f)\leq C,\quad P_{\theta,K}[\varphi]_{\mathcal{I}}(f)\leq C.
        \]
\end{corollary}
\begin{proof}
    Since $K$ is non-pluripolar, \cref{rmk:L1cpt} then gives a constant $C>0$, depending only on $\sup_K f$, $K$ and $X$, such that every $\eta\in\PSH(X,\theta)$ satisfying $\eta|_K\leq f$ quasi-everywhere satisfies $\eta\leq C$ on $X$.

    This applies to the candidates defining both envelopes. Hence both candidate families are locally uniformly bounded above, so their upper semicontinuous suprema are genuine $\theta$-psh functions. By \itemref{itm:prop:partial_env:2}, both envelopes are dominated by $f$ quasi-everywhere on $K$. Applying \cref{rmk:L1cpt} once more gives
    \[
        P_{\theta,K}[\varphi](f)\leq C,\quad P_{\theta,K}[\varphi]_{\mathcal I}(f)\leq C.
    \]
\end{proof}


\begin{proposition}\label{prop:Penv_partial_monotone}
    Let $f\colon K\rightarrow (-\infty,\infty]$ be a function bounded from below and $\varphi,\psi\in \PSH(X,\theta)_{>0}$. Suppose that $\varphi\preceq_{P} \psi$. Then
    \[
    P_{\theta,K}[\varphi](f)\leq P_{\theta,K}[\psi](f).
    \]
\end{proposition}
\begin{proof}
    Let $\eta\in \PSH(X,\theta)_{>0}$ be such that $\eta|_K\leq f$ quasi-everywhere and $\eta\preceq \varphi$. Then $\eta\preceq_P \psi$, and thanks to \cref{cor:P_par_envelope_rel} for any $C>0$, we have $\eta\land (\psi+C)\in \PSH(X,\theta)_{>0}$. But then $\eta\land (\psi+C)$ is a candidate in the definition of $P_{\theta,K}[\psi](f)$, and we find
    \[
    \eta\land (\psi+C)\leq P_{\theta,K}[\psi](f).
    \]
    By \cref{cor:P_par_envelope_rel} again, we get
    \[
    \eta=P_{\theta}[\psi](\eta)\leq P_{\theta,K}[\psi](f).
    \]
    Our assertion follows now from \cref{lma:pos_mass_partial}.
\end{proof}

\begin{corollary}\label{cor:partial_P_new_def}
    Let $f\colon K\rightarrow (-\infty,\infty]$ be a function bounded from below and $\varphi \in \PSH(X,\theta)_{>0}$. Then
    \[
    \begin{aligned}
    P_{\theta,K}[\varphi](f)= & \uscsup \Bigl\{ \eta \in \PSH(X,\theta): \eta|_K \leq f  \textup{ quasi-everywhere and } \eta\preceq_P \varphi\Bigr\}\\
    =& \uscsup \Bigl\{ \eta \in \PSH(X,\theta)_{>0}: \eta|_K \leq f  \textup{ quasi-everywhere and } \eta\preceq_P \varphi\Bigr\}\\
    =& \uscsup \Bigl\{ \eta \in \PSH(X,\theta): \eta|_K \leq f  \textup{ and } \eta\preceq_P \varphi\Bigr\}\\
    =& \uscsup \Bigl\{ \eta \in \PSH(X,\theta)_{>0}: \eta|_K \leq f  \textup{ and } \eta\preceq_P \varphi\Bigr\}.
    \end{aligned}
    \]
\end{corollary}
\begin{proof}
    We first apply \cref{prop:Penv_partial_monotone},
    \[
    \begin{split}
    P_{\theta,K}[\varphi](f)= & P_{\theta,K}\left[P_{\theta}[\varphi]\right](f)\\
    =& \uscsup \Bigl\{ \eta \in \PSH(X,\theta): \eta|_K \leq f  \textup{ quasi-everywhere and } \eta\preceq_P \varphi\Bigr\}.
    \end{split}
    \]
    Next observe that by \cref{lma:pos_mass_partial}, the latter can also be written as
    \[
    \uscsup \Bigl\{ \eta \in \PSH(X,\theta)_{>0}: \eta|_K \leq f  \textup{ quasi-everywhere and } \eta\preceq_P \varphi\Bigr\}.
    \]
    Next we apply \cref{lma:PKoutsidepps}; the remaining statements follow.
\end{proof}

\begin{corollary}\label{cor:I-good_partial_env}
    Let $f\colon K\rightarrow (-\infty,\infty]$ be a function bounded from below and $\varphi \in \PSH(X,\theta)_{>0}$ be $\mathcal{I}$-good. Then
    \[
        P_{\theta,K}[\varphi]_{\mathcal{I}}(f)=P_{\theta,K}[\varphi](f).
    \]
\end{corollary}
\begin{proof}
    In fact, since $\varphi$ is $\mathcal{I}$-good, a potential $\eta\in \PSH(X,\theta)$ satisfies $\eta\preceq_P \varphi$ if and only if $\eta\preceq_{\mathcal{I}} \varphi$. Therefore, our assertion follows from \cref{cor:partial_P_new_def} and \eqref{eq:PthetaKvarphivI}.
\end{proof}



Next, we make the following observations about the singularity types of our envelopes:
\begin{lemma}\label{lma:same_sing_type} Let $f \colon K\rightarrow \mathbb{R}$ be a bounded function, and $\varphi\in \PSH(X,\theta)_{>0}$. We have
\begin{equation}\label{eq:PthetaKvarphiv_singtype}
P_{\theta,K}[\varphi](f)\sim P_{\theta}[\varphi],\quad P_{\theta,K}[\varphi]_{\mathcal{I}}(f)\sim P_{\theta}[\varphi]_{\mathcal{I}}.
\end{equation}
\end{lemma}
\begin{proof}
    We first prove the $P$-part of \eqref{eq:PthetaKvarphiv_singtype}.

    If $\eta\in \PSH(X,\theta)$ with $\eta\preceq \varphi$ and $\eta\leq 0$, then $\eta+\inf_K f$ is a candidate in the definition \eqref{eq:PthetaKvarphiv}, and hence,
    \[
    \eta +\inf_K f \leq P_{\theta,K}[\varphi](f).
    \]
    It follows that
    \[
        P_{\theta}[\varphi]+\inf_K f \leq P_{\theta,K}[\varphi](f).
    \]

    Conversely, let $\eta\in \PSH(X,\theta)$ be such that
    \[
        \eta|_K\leq f,\quad \eta\preceq \varphi.
    \]
    Since $K$ is non-pluripolar, by \cref{rmk:L1cpt} we can find $C=C(\sup_K f,K, X)>0$ so that $\eta\leq C$. Hence $\eta-C$ is a candidate in the definition of $P_{\theta}[\varphi]$, and hence
    \[
        \eta-C\leq P_{\theta}[\varphi].
    \]
    By \cref{lma:PKoutsidepps}, this implies that
    \[
        P_{\theta,K}[\varphi](f) \leq P_{\theta}[\varphi] + C.
    \]

    The proof of the $\mathcal{I}$-part is similar, but we spell it out since this comparison will be used below. If $\eta\in \PSH(X,\theta)$ satisfies $\eta\leq 0$ and $\eta\preceq_{\mathcal{I}}\varphi$, then $\eta+\inf_K f$ is a candidate in \eqref{eq:PthetaKvarphivI}. Hence
    \[
        P_{\theta}[\varphi]_{\mathcal{I}}+\inf_K f\leq P_{\theta,K}[\varphi]_{\mathcal{I}}(f).
    \]
    Conversely, let $\eta\in \PSH(X,\theta)$ be a candidate in \eqref{eq:PthetaKvarphivI}. Since $K$ is non-pluripolar and $f$ is bounded, the same compactness estimate gives a constant $C>0$, independent of $\eta$, such that $\eta\leq C$. Then $\eta-C\leq 0$ and $\eta-C\preceq_{\mathcal{I}}\varphi$. By \cref{cor:Ienvelopedef2},
    \[
        \eta-C\leq P_{\theta}[\varphi]_{\mathcal{I}}.
    \]
    Taking the upper envelope over all such $\eta$ gives
    \[
        P_{\theta,K}[\varphi]_{\mathcal{I}}(f)\leq P_{\theta}[\varphi]_{\mathcal{I}}+C.
    \]
    This completes the proof.
\end{proof}

\begin{corollary}\label{cor:PKtoPX}
    Let $f\colon K\rightarrow \mathbb{R}$ be a bounded function, and $\varphi\in \PSH(X,\theta)_{>0}$. Then
\begin{equation}\label{eq:PthetaKvarphiveq}
P_{\theta,K}[\varphi](f)=P_{\theta,X}[\varphi]\left( P_{\theta,K}[V_{\theta}](f) \right).
\end{equation}
\end{corollary}
\begin{proof}
Take any $\eta\in \PSH(X,\theta)$ with $\eta\preceq \varphi$ and $\eta|_K\leq f$ quasi-everywhere, then
\[
\eta\leq P_{\theta,K}[\varphi](f)\leq P_{\theta,K}[V_{\theta}](f).
\]
Therefore,
\[
\eta\leq P_{\theta,X}[\varphi]\left( P_{\theta,K}[V_{\theta}](f) \right).
\]
So we conclude the $\leq$ direction in \eqref{eq:PthetaKvarphiveq}.

For the reverse direction, take $\eta\in \PSH(X,\theta)_{>0}$ such that
\[
\eta\preceq \varphi,\quad \eta\leq P_{\theta,K}[V_{\theta}](f) \quad \textup{quasi-everywhere}.
\]
Then $\eta|_K\leq f$ quasi-everywhere by \cref{prop:partial_env}. Hence
\[
\eta\leq P_{\theta,K}[\varphi](f).
\]
By \cref{cor:partial_P_new_def}, the $\geq$ direction in \eqref{eq:PthetaKvarphiveq} follows.
\end{proof}

\section{Partial equilibrium measures}\label{sec:pem}
Let $X$ be a connected compact Kähler manifold of dimension $n$, and $\theta$ be a smooth closed real $(1,1)$-form on $X$ representing a pseudo-effective cohomology class.

\begin{definition}\label{def:weightedss}
A \emph{weighted subset}\index{weighted subset} of $X$ is a pair $(K,v)$ consisting of a closed non-pluripolar subset $K\subseteq X$ and a function $v\in C^0(K)$.
\end{definition}
The next result motivates our terminology to call the measures $\theta_{P_{\theta,K}[\varphi](f)}^n$ the \emph{partial equilibrium measures}\index{partial equilibrium measure} of our context:
\begin{lemma}\label{lma:balayage}
Let $(K,v)$ be a weighted subset of $X$, and $\varphi\in \PSH(X,\theta)_{>0}$.
Then
\begin{equation}\label{eq:intXminusKthetan_vanish}
\int_{X\setminus K}\theta_{P_{\theta,K}[\varphi](v)}^n=0,\quad \int_K \mathds{1}_{\left\{P_{\theta,K}[\varphi](v)|_K\neq v\right\}}\theta_{P_{\theta,K}[\varphi](v)}^n=0.
\end{equation}
More precisely, we have
\begin{equation}\label{eq:thetaPKuv}
\theta_{P_{\theta,K}[\varphi](v)}^n\leq \mathds{1}_{\left\{P_{\theta,K}[\varphi](v)|_K=P_{\theta,K}[V_{\theta}](v)|_K=v\right\}}\,\theta_{P_{\theta,K}[V_{\theta}](v)}^n.
\end{equation}
\end{lemma}
Here $\{P_{\theta,K}[\varphi](v)|_K\neq v\}$ and $\{P_{\theta,K}[\varphi](v)|_K=P_{\theta,K}[V_{\theta}](v)|_K=v\}$ are understood as consisting of points of $K$ satisfying these constraints.


\begin{proof}
\textbf{Step~1}. We first prove \eqref{eq:intXminusKthetan_vanish} when $\varphi = V_{\theta}$.

Let $S\subseteq X$ be a proper analytic subset, such that $V_{\theta}$ is locally bounded on $X\setminus S$. Such $S$ exists by \cref{lma:Vthetalocbdd}.
Fix holomorphic balls $B\Subset B'\Subset X\setminus S$.

By \cref{prop:Choquet} and \cref{lma:PKoutsidepps}, we can take an increasing sequence $(\eta_j)_j$ in $\PSH(X,\theta)$ such that
\[
\eta_j\nearrow P_{\theta,K}[V_{\theta}](v)\textup{ a.e.},\quad \eta_j|_K\leq v\textup{ for all }j\geq 1.
\]
Note that $P_{\theta,K}[V_{\theta}](v)\sim V_{\theta}$ by \cref{lma:same_sing_type}. Take $C>0$ so that
\[
\left.\left( P_{\theta,K}[V_{\theta}](v) \right)\right|_K-C\leq v.
\]
Then we may replace $\eta_j$ by $\eta_j\lor (P_{\theta,K}[V_{\theta}](v)-C)$ and assume that $\eta_j\sim V_{\theta}$ for each $j$. In particular, $\eta_j$ is locally bounded on $X\setminus S$.

For each $j\geq 1$, let $\gamma_j\in \PSH(X,\theta)$ be the balayage of $\eta_j$ on $B$, defined in \cref{prop:bal_general}. Note that $(\gamma_j)_j$ is increasing due to \cref{prop:dom_prin_loc}.

For the first equality in \eqref{eq:intXminusKthetan_vanish},
we may assume that $B\Subset X\setminus (S \cup K)$. It suffices to show that $\theta_{P_{\theta,K}[V_{\theta}](v)}^n$ does not charge $B$.

Since $B$ lies outside $K$ by assumption, we also have $\gamma_j|_K=\eta_j|_K \leq v$ for all $j\geq 1$. In particular, $\gamma_j\leq P_{\theta,K}[V_{\theta}](v)$.

It follows that $\gamma_j \nearrow P_{\theta,K}[V_{\theta}](v)$ almost everywhere as well. Testing against non-negative continuous functions supported in $B$ and using \cref{thm:semicon}, we find that $\theta^n_{P_{\theta,K}[V_{\theta}](v)}$ does not charge $B$, as desired.

We next prove the second equality in \eqref{eq:intXminusKthetan_vanish}. The proof is essentially the same as that of \cref{lma:BT_freeenvelope}.

It suffices to show that $\theta^n_{P_{\theta,K}[V_{\theta}](v)}$ does not charge the set $K\cap \{P_{\theta,K}[V_{\theta}](v)|_K<v\}$.


Let $x\in (X\setminus S) \cap K$ be a point such that $P_{\theta,K}[V_{\theta}](v)(x)<v(x)-\epsilon$ for some $\epsilon>0$. We may assume that $x\in B$ and $B'$ is small enough, so that
\begin{enumerate}
    \item\label{itm:chapter-Bergman:L380:1} $P_{\theta,K}[V_{\theta}](v)|_{\overline {B'}}<v(x)-\epsilon$.
    \item\label{itm:chapter-Bergman:L380:2} $v|_{\overline {B'}\cap K}>v(x)-\epsilon$.
\end{enumerate}
By \ref{itm:chapter-Bergman:L380:1}, we have $\gamma_j|_{\partial B}\leq v(x)-\epsilon$, hence by \cref{prop:dom_prin_loc}, $\gamma_j|_B\leq v(x)-\epsilon$. By \ref{itm:chapter-Bergman:L380:2}, we have $\gamma_j|_{B\cap K}\leq v|_{B\cap K}$. On the other hand, $\gamma_j|_{K\setminus B}=\eta_{j}|_{K\setminus B}\leq v|_{K\setminus B}$. Therefore, $\gamma_j|_K\leq v$. It follows again that $\gamma_j\leq P_{\theta,K}[V_{\theta}](v)$. Therefore, we conclude as in the previous paragraph that $\gamma_j\nearrow P_{\theta,K}[V_{\theta}](v)$. Thus, testing against non-negative continuous functions supported in $B$, we conclude that $\theta^n_{P_{\theta,K}[V_{\theta}](v)}$ does not charge $B$ again by \cref{thm:semicon}.

\textbf{Step~2}. We handle the general case.

By \cref{cor:PKtoPX},
\[
P_{\theta,K}[\varphi](v) = P_{\theta,X}[\varphi]\left(P_{\theta,K}[V_{\theta}](v)\right).
\]
Now \cref{cor:Pvarphisupport} gives
\[
\begin{aligned}
\theta_{P_{\theta,K}[\varphi](v)}^n \leq & \mathds{1}_{\left\{P_{\theta,K}[\varphi](v) = P_{\theta,K}[V_{\theta}](v)\right\}} \theta_{P_{\theta,K}[V_{\theta}](v)}^n \\
\leq & \mathds{1}_{\left\{P_{\theta,K}[\varphi](v)|_K=P_{\theta,K}[V_{\theta}](v)|_K=v\right\}}\,\theta_{P_{\theta,K}[V_{\theta}](v)}^n,
\end{aligned}
\]
where in the second inequality we have used Step~1. Hence \eqref{eq:thetaPKuv} follows, and \emph{a fortiori} \eqref{eq:intXminusKthetan_vanish} follows as well.
\end{proof}

\begin{corollary}\label{cor:suppthetan}
Let $v\in C^0(K)$ and $\varphi\in \PSH(X,\theta)_{>0}$. Then
\begin{equation}\label{eq:thetandoesnotcharge1}
\begin{aligned}
\int_{(X\setminus K)\cup \left\{P_{\theta,K}[\varphi](v)<v\right\}}\theta^n_{P_{\theta,K}[\varphi](v)}=& 0,\\
\int_{(X\setminus K)\cup \left\{P_{\theta,K}[\varphi]_{\mathcal{I}}(v)<v\right\}}\theta^n_{P_{\theta,K}[\varphi]_{\mathcal{I}}(v)}=& 0.
\end{aligned}
\end{equation}
\end{corollary}

\begin{proof}
The first equation in \eqref{eq:thetandoesnotcharge1} follows from \cref{lma:balayage}. For the second, it suffices to observe that
\[
P_{\theta,K}[\varphi]_{\mathcal{I}}(v)=P_{\theta,K}\left[P_{\theta}[\varphi]_{\mathcal{I}}\right]_{\mathcal{I}}(v)=P_{\theta,K}\left[P_{\theta}[\varphi]_{\mathcal{I}}\right](v)
\]
by \cref{cor:I-good_partial_env}.
\end{proof}


\begin{definition}\label{def:partialequienergy}
Given $\varphi\in \PSH(X,\theta)_{>0}$, we define the \emph{partial equilibrium energy}\index{partial equilibrium energy} functional $\mathcal{E}_{\theta,K}^{\varphi}\colon C^0(K)\rightarrow \mathbb{R}$\index{$\mathcal{E}_{\theta,K}^{\varphi}$} as follows: given $v \in C^0(K)$, we set
\begin{equation}
\mathcal{E}_{\theta,K}^{\varphi}(v)\coloneqq E_{\theta}^{P_{\theta}[\varphi]_{\mathcal{I}}}\left(P_{\theta,K}[\varphi]_{\mathcal{I}}(v)\right).
\end{equation}
\end{definition}
Recall that the energy $E_{\theta}^{P_{\theta}[\varphi]_{\mathcal{I}}}$ functional is defined in \cref{def:MAenergy}.

Note that by \cref{lma:same_sing_type}, we have
\[
P_{\theta,K}[\varphi]_{\mathcal{I}}(v)\in \mathcal{E}^{\infty}\left(X,\theta;P_{\theta}[\varphi]_{\mathcal{I}}\right),
\]
so $\mathcal{E}_{\theta,K}^{\varphi}(v)\in \mathbb{R}$.


\begin{proposition}\label{prop:differential_P} Let $v,f \in C^0(K)$ and $\varphi \in \PSH(X,\theta)_{>0}$.
Then $\mathbb{R}\ni t \mapsto \mathcal{E}_{\theta,K}^{\varphi}(v+tf)$ is differentiable and
\begin{equation}\label{eq:ddtI}
\frac{\mathrm{d}}{\mathrm{d}t}\mathcal{E}_{\theta,K}^{\varphi}(v+tf) = \int_K f\, \theta^n_{P_{\theta,K}[\varphi]_{\mathcal{I}}(v+tf)}
\end{equation}
for all $t\in \mathbb{R}$.
\end{proposition}

\begin{proof}
We may assume that $\varphi$ is $\mathcal{I}$-model by replacing $\varphi$ by $P_{\theta}[\varphi]_{\mathcal{I}}$.

Note that it suffices to prove \eqref{eq:ddtI} at $t=0$, which is equivalent to
\begin{equation}\label{eq:to_prove_1}
\lim_{t\to 0}\frac{E_{\theta}^{\varphi}\left(P_{\theta,K}[\varphi]_{\mathcal{I}}(v+tf)\right)-E_{\theta}^{\varphi}\left(P_{\theta,K}[\varphi]_{\mathcal{I}}(v)\right)}{t}=\int_K f\,\theta^n_{P_{\theta,K}[\varphi]_{\mathcal{I}}(v)}.
\end{equation}
It suffices to prove the right derivative at $0$ for an arbitrary $f$; applying the same argument to $-f$ then gives the left derivative. In other words, we need to prove that
\begin{equation}\label{eq:to_prove_2}
\lim_{t\to 0+}\frac{E_{\theta}^{\varphi}\left(P_{\theta,K}[\varphi]_{\mathcal{I}}(v+tf)\right)-E_{\theta}^{\varphi}\left(P_{\theta,K}[\varphi]_{\mathcal{I}}(v)\right)}{t}=\int_K f\,\theta^n_{P_{\theta,K}[\varphi]_{\mathcal{I}}(v)}.
\end{equation}

We first reduce the proof of \eqref{eq:to_prove_2} to the case $f\leq 0$. Assume that this case is known. Choose $A\in\mathbb R$ such that $f-A\leq 0$ on $K$. Since for any $t>0$,
\[
P_{\theta,K}[\varphi]_{\mathcal I}\Bigl( v+t(f-A) \Bigr)= P_{\theta,K}[\varphi]_{\mathcal I}(v+tf)-tA,
\]
then we have
\[
    \int_K f\, \theta^n_{P_{\theta,K}[\varphi]_{\mathcal{I}}(v)}=\int_K (f-A)\, \theta^n_{P_{\theta,K}[\varphi]_{\mathcal{I}}(v)}+A\int_X \theta^n_{P_{\theta,K}[\varphi]_{\mathcal{I}}(v)}.
\]
By \eqref{eq:Ephiaddcst}, for any $t>0$,
\[
    E_{\theta}^{\varphi}\Bigl(P_{\theta,K}[\varphi]_{\mathcal{I}}\left(v+t(f-A)\right)\Bigr)=E_{\theta}^{\varphi}\left(P_{\theta,K}[\varphi]_{\mathcal{I}}(v+tf)\right)-tA\int_X\theta_{\varphi}^n.
\]
So the limit in \eqref{eq:to_prove_2} exists and is equal to the left-hand side of \eqref{eq:to_prove_2} with $f-A$ in place of $f$ plus $A\int_X\theta_{\varphi}^n$.
So \eqref{eq:to_prove_2} for $f-A$ implies the desired formula for $f$ in view of \cref{lma:same_sing_type} and the monotonicity theorem \cref{thm:mono}. Hence we may assume that $f\leq 0$.


Fix $t>0$.
Set
\[
\eta_t\coloneqq P_{\theta,K}[\varphi]_{\mathcal I}(v+tf),\quad \eta_0\coloneqq P_{\theta,K}[\varphi]_{\mathcal I}(v).
\]
By the comparison principle \eqref{eq:comp_princi1} and \cref{prop:cocycE1}, we find
\[
\begin{aligned}
E_{\theta}^{\varphi}(\eta_t)-E_{\theta}^{\varphi}(\eta_0)
={}&\frac{1}{n+1}\sum_{i=0}^n\int_X(\eta_t-\eta_0)\,
\theta_{\eta_t}^i\wedge\theta_{\eta_0}^{n-i}\\
\leq{}& \int_X(\eta_t-\eta_0)\,\theta_{\eta_0}^n.
\end{aligned}
\]
By \cref{lma:balayage} and \cref{cor:I-good_partial_env}, the measure $\theta_{\eta_0}^n$ is carried by $K\cap\{\eta_0=v\}$, and the measure $\theta_{\eta_t}^n$ is carried by $K\cap\{\eta_t=v+tf\}$.
On $K\cap\{\eta_0=v\}$, we have $\eta_t\leq v+tf=\eta_0+tf$, hence
\[
\eta_t-\eta_0\leq tf.
\]
On $K\cap\{\eta_t=v+tf\}$, we have $\eta_0\leq v$, hence
\[
\eta_t-\eta_0\geq tf.
\]
It follows that
\[
\int_X(\eta_t-\eta_0)\,\theta_{\eta_0}^n \leq t\int_K f\,\theta_{\eta_0}^n
\]
and
\[
\int_X(\eta_t-\eta_0)\,\theta_{\eta_t}^n \geq t\int_K f\,\theta_{\eta_t}^n.
\]
Thus, letting $t\to 0+$, we get the inequality
\[
\varlimsup_{t\to 0+}\frac{E_{\theta}^{\varphi}\left(P_{\theta,K}[\varphi]_{\mathcal{I}}(v+tf)\right)-E_{\theta}^{\varphi}\left(P_{\theta,K}[\varphi]_{\mathcal{I}}(v)\right)}{t}\leq\int_K f\,\theta^n_{P_{\theta,K}[\varphi]_{\mathcal{I}}(v)}.
\]
Similarly, for $t>0$, we have
\[
\begin{aligned}
&E_{\theta}^{\varphi}\left(P_{\theta,K}[\varphi]_{\mathcal{I}}(v+tf)\right)-E_{\theta}^{\varphi}\left(P_{\theta,K}[\varphi]_{\mathcal{I}}(v)\right)\\
\geq & \int_X \left(P_{\theta,K}[\varphi]_{\mathcal{I}}(v+tf)-P_{\theta,K}[\varphi]_{\mathcal{I}}(v)\right)\,\theta_{P_{\theta,K}[\varphi]_{\mathcal{I}}(v+tf)}^n\\
\geq & t\int_K f\,\theta_{P_{\theta,K}[\varphi]_{\mathcal{I}}(v+tf)}^n.
\end{aligned}
\]
Letting $t\to 0+$, we find
\[
\varliminf_{t\to 0+}\frac{E_{\theta}^{\varphi}\left(P_{\theta,K}[\varphi]_{\mathcal{I}}(v+tf)\right)-E_{\theta}^{\varphi}\left(P_{\theta,K}[\varphi]_{\mathcal{I}}(v)\right)}{t}\geq\int_K f\,\theta^n_{P_{\theta,K}[\varphi]_{\mathcal{I}}(v)}.
\]
Here we used the following elementary continuity consequence:
\[
P_{\theta,K}[\varphi]_{\mathcal{I}}(v)-t\|f\|_{C^0(K)}
\leq P_{\theta,K}[\varphi]_{\mathcal{I}}(v+tf)
\leq P_{\theta,K}[\varphi]_{\mathcal{I}}(v)
\]
for $t>0$, and hence $P_{\theta,K}[\varphi]_{\mathcal{I}}(v+tf) \nearrow P_{\theta,K}[\varphi]_{\mathcal{I}}(v)$ almost everywhere when $t\to 0+$.
Next we claim that
\[
    \int_K f\,\theta^n_{P_{\theta,K}[\varphi]_{\mathcal{I}}(v)}\leq \lim_{t\to 0+}  \int_K f\,\theta_{P_{\theta,K}[\varphi]_{\mathcal{I}}(v+tf)}^n.
\]
Choose a continuous extension $\tilde f\in C^0(X)$ of $f$. Since \cref{lma:balayage,cor:I-good_partial_env} show that all the measures appearing here are carried by $K$, the integrals over $K$ against $f$ are the same as the integrals over $X$ against $\tilde f$. Thus \cref{cor:incseqnppcont} applies.

We conclude \eqref{eq:to_prove_1}.
\end{proof}

\begin{lemma}\label{lma:global_env_approx} 
    Assume that $\{\theta\}$ represents a big cohomology class.
    Fix a Kähler form $\omega$ on $X$ and a weighted subset $(K,v)$ of $X$. There exists a uniformly bounded increasing sequence $(v^-_j)_j$ in $C^\infty(X)$ and a uniformly bounded decreasing sequence $(v^+_j)_j$ in $C^\infty(X)$, such that
    \begin{equation}\label{eq:unf_bd_vj}
    \sup_X |v_j^-|\leq C,\quad \sup_X |v_j^+|\leq C
    \end{equation}
    for some constant $C$ depending only on $\|v\|_{C^0(K)}$, $X$, $K$ and $\theta+\omega$. Furthermore, for any $\varphi \in \PSH(X,\theta)_{>0}$ and any $\delta\in [0,1]$ we have the following:
\begin{enumerate}
    \item\label{itm:lma:global_env_approx:1} We have
    \[
    P_{\theta+\delta \omega,X}[\varphi](v_j^+) \searrow P_{\theta+\delta \omega,K}[\varphi](v),
    \]
    and
    \[
    P_{\theta+\delta \omega,X}[\varphi](v^-_j) \nearrow P_{\theta+\delta \omega,K}[\varphi](v)
    \]
    almost everywhere.
    \item\label{itm:lma:global_env_approx:2} Suppose in addition that $\varphi$ is $\mathcal{I}$-good, then
    \[
        \lim_{j\to\infty}\mathcal{E}_{\theta,X}^{\varphi}(v_j^-)=\mathcal{E}_{\theta,K}^{\varphi}(v),\quad \lim_{j\to\infty}\mathcal{E}_{\theta,X}^{\varphi}(v_j^+)=\mathcal{E}_{\theta,K}^{\varphi}(v).
    \]
\end{enumerate}
\end{lemma}
\begin{proof}
    \textbf{Step~1}. First we construct $(v^-_j)_j$.

    Let
    \[
        C_{K,v} \coloneqq \sup\left\{\sup_X \eta : \eta \in \PSH(X,\theta+\omega),  \eta|_K \leq v\right\}.
    \]
    Since $K$ is non-pluripolar, we have $C_{K,v} \in \mathbb{R}$ by \cref{rmk:L1cpt}.
    Set
    \[
    M_{K,v}\coloneqq \max\{C_{K,v}+1,\sup_K v\}.
    \]
    Now define $\widetilde v \colon X \to \mathbb{R}$ as
    \[
    \widetilde{v}(x)=\left\{
    \begin{aligned}
        v(x),&\quad x\in K;\\
        M_{K,v},&\quad x\in X\setminus K.
    \end{aligned}
    \right.
    \]
    Since $M_{K,v}\geq \sup_K v$, the function $\widetilde v$ is lower semicontinuous. Therefore there exists an increasing and uniformly bounded sequence $(v^-_j)_j$ in $C^\infty(X)$ such that $v^-_j \nearrow \widetilde v$ and
    \[
        \|v_j^-\|_{C^0(X)}\leq \|\tilde{v}\|_{L^{\infty}(X)}+1.
    \]
    The right-hand side can clearly be bounded by a constant as in \eqref{eq:unf_bd_vj}.

    \textbf{Step~2}. We construct $(v_j^+)_j$.

    Since
\[
h \coloneqq P_{\theta + \omega,K}[V_{\theta+\omega}](v) \lor \left(\inf_K v -1\right)
\]
is usc, there exists a decreasing and uniformly bounded sequence $(v^+_j)_j$ in $C^\infty(X)$ such that $v^+_j \searrow h$. We may guarantee that
\[
\|v^+_j\|_{C^0(X)}\leq \|h\|_{L^{\infty}(X)}+1.
\]
Note that $\inf_K v-1\leq h\leq C_{K,v} \lor (\inf_K v-1)$. So \eqref{eq:unf_bd_vj} is satisfied.




\textbf{Step~3}. We finish the proof. Fix $\varphi\in \PSH(X,\theta)_{>0}$ and $\delta\in [0,1]$.

Note that \ref{itm:lma:global_env_approx:2} follows from \cref{cor:I-good_partial_env}, \cref{lma:same_sing_type}, \cref{prop:Econt_mono} and \ref{itm:lma:global_env_approx:1}. Let us prove \ref{itm:lma:global_env_approx:1}.

\textbf{Step~3.1}. We prove \ref{itm:lma:global_env_approx:1} for $(v^-_j)_j$.

Observe that $P_{\theta + \delta \omega,X}[\varphi](v^-_j)$ is increasing in $j\geq 1$, and
\[
P_{\theta + \delta \omega,X}[\varphi]\left(v^-_j\right)\leq P_{\theta + \delta \omega,K}[\varphi](v).
\]
We claim that
\begin{equation}\label{eq:Pthetapdeltaomega_temp1}
P_{\theta + \delta \omega,X}[\varphi]\left(v^-_j\right) \nearrow P_{\theta + \delta \omega,K}[\varphi](v)\quad \textup{a.e.}.
\end{equation}
Let $\eta\in \PSH(X,\theta+\delta\omega)$ be a potential such that $\eta\preceq \varphi$ and $\eta|_K\leq v$. Take some $\epsilon>0$. Then, since $\eta$ is upper semicontinuous, $\widetilde v$ is lower semicontinuous and $\eta < \epsilon+\widetilde v$, compactness gives $j_0>0$ such that $\eta < \epsilon+ v^-_j$ for $j \geq j_0$. Hence,
\[
\eta \leq \epsilon+P_{\theta + \delta \omega,X}[\varphi](v^-_j)
\]
for $j\geq j_0$.
Therefore,
\[
\eta\leq \uscsup[j\geq 1] P_{\theta + \delta \omega,X}[\varphi]\left(v^-_j\right)+\epsilon.
\]
Letting $\epsilon\to 0+$, we find
\[
\eta\leq \uscsup[j\geq 1] P_{\theta + \delta \omega,X}[\varphi]\left(v^-_j\right).
\]
Since $\eta$ is arbitrary, \eqref{eq:Pthetapdeltaomega_temp1} follows.

\textbf{Step~3.2}. We prove \ref{itm:lma:global_env_approx:1} for $(v^+_j)_j$.

Note that $P_{\theta + \delta \omega,X}[\varphi](v^+_j)$ is decreasing in $j$, and
\begin{equation}\label{eq:defchilimPthetadeltaomega}
\chi \coloneqq \lim_{j\to\infty} P_{\theta + \delta \omega,X}[\varphi]\left(v^+_j\right) \geq P_{\theta + \delta \omega,K}[\varphi](v).
\end{equation}
In particular, $\chi\in \PSH(X,\theta+\delta\omega)_{>0}$, since $P_{\theta + \delta \omega,K}[\varphi](v)$ has the same singularity types as $P_{\theta + \delta \omega}[\varphi]$ by \cref{lma:same_sing_type}.
Observe that
\[
P_{\theta + \delta \omega,X}[\varphi]\left(v_j^+\right)\leq v_j^+
\]
for all $j\geq 1$, since $v_j^+$ is continuous. Thus, $\chi\leq h$.
On the other hand, $h|_K\leq v$ quasi-everywhere. So in particular, $\chi|_{K}\leq v$ quasi-everywhere.
On the other hand,
\[
\chi\leq P_{\theta + \delta \omega,X}[\varphi]\left(v^+_1\right)\sim_P \varphi
\]
by \cref{lma:same_sing_type}. Hence, by \cref{cor:partial_P_new_def},
\[
\chi\leq P_{\theta + \delta \omega,K}[\varphi](v).
\]
Therefore, \eqref{eq:defchilimPthetadeltaomega} is an equality.
\end{proof}


\begin{proposition}\label{prop:conv_of_K_env}
    Let $(K,v)$ be a weighted subset of $X$. Let $(\varphi_j)_j$ be a sequence in $\PSH(X,\theta)_{>0}$ and $\varphi\in \PSH(X,\theta)_{>0}$. Then the following hold:
\begin{enumerate}
    \item\label{itm:prop:conv_of_K_env:1} If $(\varphi_j)_j$ is decreasing, and $\varphi_j\xrightarrow{d_S}\varphi$, then
    \begin{equation}\label{eq:PthetaKvarphij1}
    P_{\theta,K}[\varphi_j]_{\mathcal{I}}(v) \searrow P_{\theta,K}[\varphi]_{\mathcal{I}}(v),\quad P_{\theta,K}[\varphi_j](v) \searrow P_{\theta,K}[\varphi](v).
    \end{equation}
    \item\label{itm:prop:conv_of_K_env:2} If $\varphi_j \nearrow \varphi$ almost everywhere then
    \[
    P_{\theta,K}[\varphi_j]_{\mathcal{I}}(v) \nearrow P_{\theta,K}[\varphi]_{\mathcal{I}}(v)\textup{ a.e.},\quad P_{\theta,K}[\varphi_j](v) \nearrow P_{\theta,K}[\varphi](v) \textup{ a.e.}.
    \]
\end{enumerate}
\end{proposition}

\begin{proof}
We only give the proof for the relative $\mathcal{I}$-envelopes; the proof for the relative $P$-envelopes is identical, replacing $P_{\theta}[\bullet]_{\mathcal{I}}$ by $P_{\theta}[\bullet]$ and using \cref{cor:partial_P_new_def} for the final candidate condition.

Set
\[
    \eta_j\coloneqq P_{\theta,K}[\varphi_j]_{\mathcal{I}}(v),\quad
    \eta\coloneqq P_{\theta,K}[\varphi]_{\mathcal{I}}(v).
\]
By \cref{lma:same_sing_type}, $d_S(\eta_j,P_{\theta}[\varphi_j]_{\mathcal{I}})=0$ and $d_S(\eta,P_{\theta}[\varphi]_{\mathcal{I}})=0$ for all $j$. Moreover, \cref{thm:contPI} gives
\[
    P_{\theta}[\varphi_j]_{\mathcal{I}}\xrightarrow{d_S}P_{\theta}[\varphi]_{\mathcal{I}}
\]
in case \ref{itm:prop:conv_of_K_env:1}, and the same convergence holds in case \ref{itm:prop:conv_of_K_env:2} by \cref{cor:incnetdSconv} and \cref{thm:contPI}.
Consequently, $\eta_j\xrightarrow{d_S}\eta$ in both cases.

Assume first that $(\varphi_j)_j$ is decreasing. Then $(\eta_j)_j$ is decreasing and $\eta_j\geq \eta$. Let $\eta_\infty\coloneqq \inf_j \eta_j$. The Monge--Amp\`ere masses of $\eta_j$ converge to the mass of $\eta$, by the preceding $d_S$-convergence and \cref{thm:convdS}. Since $\eta_j\geq \eta_\infty\geq \eta$, the monotonicity theorem \cref{thm:mono} gives
\[
    \int_X\theta_{\eta_\infty}^n=\int_X\theta_\eta^n.
\]
Hence $\eta_j\xrightarrow{d_S}\eta_\infty$ by \cref{cor:decnetdS}. The uniqueness of the $d_S$-limit gives $d_S(\eta_\infty,\eta)=0$. Since $\eta_\infty\geq \eta$, \cref{prop:ds0char} implies $\eta_\infty\sim_P \eta$. In particular, $\eta_\infty\preceq_{\mathcal{I}}\varphi$ by \cref{cor:PimpliesI}. The condition $\eta_\infty\leq v$ on $K$ quasi-everywhere is inherited from the $\eta_j$'s, so $\eta_\infty$ is a candidate for $P_{\theta,K}[\varphi]_{\mathcal{I}}(v)$. Thus $\eta_\infty\leq \eta$, and the reverse inequality was already observed. This proves the $\mathcal{I}$-part of \ref{itm:prop:conv_of_K_env:1}.

Assume next that $\varphi_j\nearrow\varphi$ almost everywhere. Then $(\eta_j)_j$ is increasing and $\eta_j\leq \eta$. Put $\tilde\eta\coloneqq \uscsup_j \eta_j$. By \cref{cor:Ppartbdd_bdabove} and \cref{cor:incnetdSconv}, we have $\eta_j\xrightarrow{d_S}\tilde\eta$. The $d_S$-convergence above therefore gives $d_S(\tilde\eta,\eta)=0$, hence $\tilde\eta\sim_P \eta$ by \cref{prop:ds0char}. In particular, $\tilde\eta\preceq_{\mathcal{I}}\varphi$ by \cref{cor:PimpliesI}. Also $\tilde\eta\leq v$ on $K$ quasi-everywhere by \cref{prop:supsupstardiff}. Hence $\tilde\eta$ is a candidate for $P_{\theta,K}[\varphi]_{\mathcal{I}}(v)$, so $\tilde\eta\leq \eta$. Since $\eta_j\leq \eta$ for all $j$, this proves $\eta_j\nearrow \eta$ almost everywhere.
\end{proof}







\section{Quantization of partial equilibrium measures}\label{sec:quant}


Let $X$ be a connected compact K\"ahler manifold of dimension $n$, let $L$ be a pseudo-effective line bundle on $X$, and let $h$ be a Hermitian metric on $L$. Set $\theta=c_1(L,h)$. Let $(T,h_T)$ be a Hermitian line bundle on $X$.

\subsection{Bernstein--Markov measures}
\label{subsec:BM}
Let $K\subseteq X$ be a closed non-pluripolar subset, and $v\colon K\rightarrow [-\infty,\infty)$ be a measurable function.

\begin{definition}\label{def:Bergman:L706}
    Let $\mu$ be a Borel probability measure on $K$.
We introduce the following functions on $\mathrm{H}^0(X,L^k \otimes T)$ ($k\geq 1$), with values possibly equal to $\infty$:
\index{$N^k_{v,\mu}$}
\index{$N^k_{v,K}$}
\begin{equation}\label{eq:Nkvmu_norm}
\begin{aligned}
N^k_{v,\mu}(s) \coloneqq & \left(\int_K h^k \otimes h_T(s,s) \mathrm{e}^{-kv} \,\mathrm{d}\mu\right)^{1/2},\\
N^k_{v,K}(s) \coloneqq & \sup_{K\setminus \{v=-\infty\}} \left( h^k \otimes h_T(s,s)\mathrm{e}^{-kv}\right)^{1/2}.
\end{aligned}
\end{equation}
\end{definition}
Observe that if $\mu$ puts no mass on $\{v=-\infty\}$ then we always have
\begin{equation}\label{eq:Nkinfcomp}
N^k_{v,\mu}(s)\leq N^k_{v,K}(s).
\end{equation}


We start with the following elementary observation:
\begin{lemma}\label{lma:mononorm}
Let $v_1,v_2\colon K\rightarrow [-\infty,\infty)$ be two measurable functions such that
\begin{enumerate}
    \item\label{itm:lma:mononorm:1} $v_1\leq v_2$;
    \item\label{itm:lma:mononorm:2} $\{v_1=-\infty\}=\{v_2=-\infty\}$.
\end{enumerate}
Then for any $s\in \mathrm{H}^0(X,L^k\otimes T)$ ($k\geq 1$), we have
\[
N_{v_1,K}^k(s)\geq N_{v_2,K}^k(s).
\]
\end{lemma}
\begin{proof}
    This follows immediately from the definition \eqref{eq:Nkvmu_norm}.
\end{proof}






\begin{definition}\label{def:BMmeasure}
Let $(K,v)$ be a weighted subset of $X$.
A Borel probability measure $\mu$ on $K$ is \emph{Bernstein--Markov}\index{Bernstein--Markov measure} with respect to $(K,v)$ if for each $\epsilon>0$, there is a constant $C_\epsilon>0$ such that
\begin{equation}\label{eq:BM}
     N^k_{v,K}(s) \leq C_\epsilon \mathrm{e}^{\epsilon k} N^k_{v,\mu}(s)
\end{equation}
for any $s \in \mathrm{H}^0(X,L^k \otimes T)$ and any $k\in \mathbb{Z}_{>0}$.
We write $\BM(K,v)$\index{$\BM(K,v)$} for the set of Bernstein--Markov measures with respect to $(K,v)$.
\end{definition}
As pointed out in \cite{BBWN11}, any volume form on $X$ is Bernstein--Markov with respect to $(X,v)$, with $v \in C^\infty(X)$.


\begin{proposition}\label{prop:BMNimplynorm}
Assume that $(K,v)$ is a weighted subset of $X$. Then
\begin{enumerate}
    \item\label{itm:prop:BMNimplynorm:1} $N^k_{v,K}$ is a norm on $\mathrm{H}^0(X,L^k\otimes T)$.
    \item\label{itm:prop:BMNimplynorm:2} For any $\mu\in \BM(K,v)$, $N^k_{v,\mu}$ is a norm on $\mathrm{H}^0(X,L^k\otimes T)$.
\end{enumerate}
\end{proposition}
\begin{proof}
\ref{itm:prop:BMNimplynorm:1}
As $v$ is bounded, $N^k_{v,K}$ is clearly finite on $\mathrm{H}^0(X,L^k\otimes T)$. In order to show that it is a norm, it suffices to show that for any $s\in \mathrm{H}^0(X,L^k\otimes T)$, $N^k_{v,K}(s)=0$ implies that $s=0$. In fact, we have $s|_K=0$, hence $s=0$ by the connectedness of $X$ and the assumption that $K$ is non-pluripolar.

\ref{itm:prop:BMNimplynorm:2}
As $v$ is bounded, clearly $N^k_{v,\mu}$ is finite and satisfies the triangle inequality. Non-degeneracy follows from the fact that $N^k_{v,K}$ is a norm and \eqref{eq:BM}.
\end{proof}



\subsection{Partial Bergman kernels}
\label{subsec:pBerg}
In this section, we fix a weighted subset $(K,v)$ of $X$ and $\mu\in \BM(K,v)$.

\begin{definition}\label{def:Bergman:L778}
    For any $\varphi\in \PSH(X,\theta)_{>0}$, we introduce the \emph{partial Bergman kernels}\index{partial Bergman kernels} of $\varphi$ (with respect to $(K,v)$) as follows: For any $k\geq 1$, we introduce
    \begin{equation}
    \begin{aligned}
     B^k_{v,\varphi, \mu}(x) \coloneqq \sup \left\{h^k \otimes h_T(s,s)\mathrm{e}^{-kv}(x):s \in \mathrm{H}^0\left(X,L^k \otimes T \otimes \mathcal{I}(k\varphi)\right),\right.\\
     \left. N^k_{v,\mu}(s) \leq 1 \right\},\quad x\in K.
     \end{aligned}
 \end{equation}
 \index{$B^k_{v,\varphi, \mu}$}
 We extend $B^k_{v,\varphi, \mu}$ to the whole $X$ by setting it to be $0$ outside $K$.

 The \emph{partial Bergman measures}\index{partial Bergman measures} of $\varphi$ (with respect to $(K,v)$) are defined as
 \begin{equation}
\beta^k_{v,\varphi,\mu} \coloneqq \frac{n!}{k^n} B^k_{v,\varphi,\mu} \,\mathrm{d}\mu
\end{equation}
\index{$\beta^k_{v,\varphi,\mu}$}
for each $k\geq 1$.
\end{definition}

\begin{remark}\label{rmk:Bergman:L797}
    Recall that in the definition of $\beta^k_{v,\varphi,\mu}$, the singular metric $\varphi$ only appears through the conditions $s \in \mathrm{H}^0(X,L^k \otimes T \otimes \mathcal{I}(k\varphi))$. The pointwise norm used to define $\beta^k_{v,\varphi,\mu}$ and $N^k_{v,\mu}$ is defined via $v$ not $\varphi$.

    This leads to the following seemingly paradoxical consequence, if $\psi\in \PSH(X,\theta)$ and $\varphi\preceq_{\mathcal{I}}\psi$, then
    \[
        \beta^k_{v,\varphi,\mu}\leq \beta^k_{v,\psi,\mu}\quad \forall\,k\geq 1.
    \]
    We shall apply this consequence without further mention in the sequel.
\end{remark}


Observe that
\[
    \int_X \beta^k_{v,\varphi,\mu}=\frac{n!}{k^n}h^0\left(X,T\otimes L^k \otimes \mathcal{I}(k\varphi)\right).
\]
In particular,
\begin{equation}\label{eq:betaint_conv}
    \lim_{k\to\infty}\int_X \beta^k_{v,\varphi,\mu}=\vol\left(\theta_{\varphi}\right)=\int_X \theta_{P_{\theta,K}[\varphi]_{\mathcal{I}}(v)}^n
\end{equation}
by \cref{thm:DXmain1} and \cref{lma:same_sing_type}. The goal of this section is to improve this result to the convergence of the Bergman kernels themselves:
\begin{theorem}\label{thm:pBMconvergence} Suppose that $\varphi \in \PSH(X,\theta)_{>0}$. Then
\begin{equation}\label{eq:pbkconvgeneral}
\beta^k_{v,\varphi,\mu} \rightharpoonup \theta_{P_{\theta,K}[\varphi]_{\mathcal{I}}(v)}^n,\quad \textup{as }k\to\infty.
\end{equation}
\end{theorem}

The proof will be given at the end of this subsection. We shall prove a few preliminaries at first.



Take a K\"ahler form $\omega$ on $X$ so that
\[
\int_X \omega^n = 1.
\]
\begin{proposition}\label{prop:smooth_weak_conv}
    Assume here that $K=X$.
    Let $\varphi \in \PSH(X,\theta)$ be a potential with analytic singularities such that $\theta_{\varphi}$ is a Kähler current. If $v \in C^\infty(X)$, then
\begin{equation}\label{eq:pbmconvanaly}
\beta^k_{v,\varphi,\omega^n} \rightharpoonup \theta_{P_{\theta,X}[\varphi]_{\mathcal{I}}(v)}^n=\theta_{P_{\theta,X}[\varphi](v)}^n\quad \textup{as }k\to\infty.
\end{equation}
\end{proposition}

\begin{proof} The equality part in \eqref{eq:pbmconvanaly} follows from \cref{cor:I-good_partial_env}. It remains to prove the weak convergence.

Let $\tau$ be the weak limit of a subsequence of $(\beta^k_{v,\varphi,\omega^n})_k$. In view of \eqref{eq:betaint_conv}, it suffices to show that
\begin{equation}\label{eq:muleq_temp1}
\tau\leq \theta_{P_{\theta,X}[\varphi](v)}^n.
\end{equation}

We can obtain an upper bound of $\tau$ via the observation:
\[
\beta^k_{v,\varphi,\omega^n} \leq \beta^k_{v,V_{\theta},\omega^n}.
\]
After replacing $\varphi$ by $V_{\theta}$, the corresponding problem is handled by \cite[Theorem~1.2]{Ber11}, giving us
\[
\beta^k_{v,V_{\theta},\omega^n} \rightharpoonup \theta_{P_{\theta,X}[V_{\theta}](v)}^n=\mathds{1}_{\left\{v = P_{\theta,X}[V_{\theta}](v)\right\}}\,\theta_v^n \quad \textup{as }k\to\infty,
\]
where the last identity is due to \cref{prop:DNTenvelope2}.

We obtain that
\begin{equation}\label{eq:Bergmanmeasure}
\tau \leq \mathds{1}_{\left\{v = P_{\theta,X}[V_{\theta}](v)\right\}}\,\theta_v^n.
\end{equation}
Now observe that it suffices to show that
\begin{equation}\label{eq:muP0_temp1}
\tau\left\{P_{\theta,X}[\varphi](v) < v\right\}=0.
\end{equation}
Assuming this, then \eqref{eq:Bergmanmeasure} further gives
\[
\tau \leq \mathds{1}_{\left\{P_{\theta,X}[\varphi](v)= v\right\}}\,\theta_v^n= \theta_{P_{\theta,X}[\varphi](v)}^n,
\]
where the equality follows from \cref{prop:DNTenvelope2}, and \eqref{eq:muleq_temp1} follows.

It now remains to argue \eqref{eq:muP0_temp1}. Fix a Kähler form $\omega$ on $X$. Observe that
\[
P_{\theta +\epsilon \omega,X}[\varphi](v) \sim \varphi
\]
for all $\epsilon\geq 0$ by \cref{lma:same_sing_type} and \cref{prop:analysingcompPandPI}. The family $P_{\theta+\epsilon\omega,X}[\varphi](v)$ is decreasing as $\epsilon\to0+$. Let
\[
    \chi\coloneqq \inf_{\epsilon>0}P_{\theta+\epsilon\omega,X}[\varphi](v).
\]
Then $\chi\in\PSH(X,\theta)$ since $\chi\geq P_{\theta,X}[\varphi](v)$. Moreover $\chi\leq v$ and $\chi\preceq\varphi$. Hence
\[
    \chi\leq P_{\theta,X}[\varphi](v).
\]
Thus, equality holds, and
\[
    P_{\theta+\epsilon\omega,X}[\varphi](v)\searrow P_{\theta,X}[\varphi](v).
\]
In other words,
\[
\left\{P_{\theta,X}[\varphi](v) < v\right\}=\bigcup_{\epsilon\in \mathbb{Q}_{>0}}\left\{P_{\theta+\epsilon\omega,X}[\varphi](v) < v\right\}.
\]
So in order to show \eqref{eq:muP0_temp1}, we may fix $\epsilon>0$ and argue that
\begin{equation}\label{eq:muP0_temp2}
\tau\left\{P_{\theta+\epsilon\omega,X}[\varphi](v) < v\right\}=0.
\end{equation}



Let $k\geq 1$, and $s \in \mathrm{H}^0(X,L^k \otimes T \otimes \mathcal I(k\varphi))$ be a section such that $N^k_{v,\omega^n}(s) \leq 1$. Then by \cite[Lemma~4.1]{Ber11}, there exists $C>0$ such that
\[
h^k\otimes h_T(s,s)\mathrm{e}^{-kv} \leq B^k_{v,\varphi,\omega^n} \leq B^k_{v,V_{\theta},\omega^n} \leq k^n C.
\]
This implies that
\[
\frac{1}{k}\log h^k \otimes h_T(s,s) \leq v + \frac{\log C }{k} + n \frac{\log k}{ k }.
\]
We define $\varphi_k$ as in \cref{prop:Bergman_approx}. Take $\alpha_k\nearrow 1$ as in \cref{prop:Bergman_approx}.
After adding a constant to $\varphi$, we may assume that $\varphi\leq 0$ on $X$. After replacing $\alpha_k$ by a subsequence, we also have $\alpha_k\varphi\in \PSH(X,\theta+\epsilon\omega)$ for all $k>0$, and $\alpha_k\varphi\searrow\varphi$.
Then
\[
\frac{1}{k}\log h^k \otimes h_T(s,s)\preceq \varphi_k\preceq \alpha_k\varphi.
\]
We notice that $\frac{1}{k}\log h^k \otimes h_T(s,s) \in \PSH(X, \theta + \epsilon \omega)$ for all $k \geq k_0(\epsilon)$.
Indeed,
\[
    \frac{1}{k}\log h^k\otimes h_T(s,s) \in \PSH\left(X,\theta+ k^{-1} c_1(T,h_T)\right)
\]
by \cref{prop:LelongPoincare}.


In particular,
\[
\frac{1}{k}\log h^k \otimes h_T(s,s)-\frac{\log C}{k}-n\frac{\log k}{k}\leq P_{\theta+\epsilon\omega,X}[\alpha_k \varphi](v).
\]
Now taking supremum over all candidates $s$, we obtain that
\begin{equation}\label{eq:smooth_Berg_est}
B^k_{v,\varphi,\omega^n} \leq C k^n \mathrm{e}^{k \left(P_{\theta +\epsilon \omega,X}[\alpha_k \varphi](v)-v\right)},\quad k \geq k_0.
\end{equation}

Now we prove \eqref{eq:muP0_temp2}. First observe that $\alpha_k\varphi\xrightarrow{d_S}\varphi$ by \cref{lma:dSsmallmult}. Thus, by \cref{prop:conv_of_K_env},
\[
P_{\theta +\epsilon \omega,X}[\alpha_k \varphi](v) \searrow P_{\theta +\epsilon \omega,X}[\varphi](v)
\]
and we get that
\[
\bigcup_{k=1}^{\infty}\left\{P_{\theta +\epsilon \omega,X}[\alpha_k \varphi](v) < v\right\}=\left\{P_{\theta +\epsilon \omega,X}[\varphi](v) < v\right\}.
\]
    It suffices to show that $\tau$ does not put mass on the set $\{P_{\theta +\epsilon \omega,X}[\alpha_m \varphi](v) < v\}$ for each fixed $m$. Fix a positive rational $\delta>0$, it suffices to show that $\tau$ puts no mass on $F\coloneqq \{P_{\theta +\epsilon \omega,X}[\alpha_m \varphi](v) < v-\delta\}$.
    Note that $F$ is open, and \eqref{eq:smooth_Berg_est} gives $\beta^\ell_{v,\varphi,\omega^n}(F)\to 0$. Therefore, $\tau$ puts no mass on $F$ as well, proving \eqref{eq:muP0_temp2}.
\end{proof}



\begin{definition}\label{def:Bergman:L942}
    Take $k\geq 1$ and $\varphi\in \PSH(X,\theta)$, let $\Nm \mathrm{H}^0(X,L^k \otimes T \otimes \mathcal{I}(k\varphi))$ be the space of norms on the vector space $\mathrm{H}^0(X,L^k \otimes T \otimes \mathcal{I}(k\varphi))$.

    Let $\mathcal{L}_{k,\varphi} \colon \Nm \mathrm{H}^0(X,L^k \otimes T \otimes \mathcal{I}(k\varphi)) \to \mathbb R$ be the \emph{partial Donaldson functional}\index{partial Donaldson functional}:
\begin{equation}
\mathcal{L}_{k,\varphi}(H) = \frac{n!}{k^{n+1}} \log \frac{\vol\{s \in \mathrm{H}^0\left(X,L^k \otimes T \otimes \mathcal{I}(k\varphi)\right) : H(s) \leq 1\}}{\vol \left\{s\in \mathrm{H}^0\left(X,L^k \otimes T \otimes \mathcal{I}(k\varphi)\right) : N^k_{0,\omega^n}(s) \leq 1\right\}},
\end{equation}
\index{$\mathcal{L}_{k,\varphi}$}
where $\vol$ is a Lebesgue measure on $\mathrm{H}^0(X,L^k \otimes T \otimes \mathcal{I}(k\varphi))$. The normalization does not matter.
\end{definition}

\begin{proposition}\label{prop:quant_I_smooth}
Let $w,w' \in C^0(X)$ and $\varphi \in \PSH(X,\theta)$ be a potential with analytic singularities such that $\theta_{\varphi}$ is a Kähler current. Then
\begin{equation}\label{eq:LdiffonXsmoothmeasure}
\lim_{k\to\infty} \left(\mathcal{L}_{k,\varphi}(N^k_{w,\omega^n}) - \mathcal{L}_{k,\varphi}(N^k_{w',\omega^n}) \right)= \mathcal{E}_{\theta,X}^{\varphi}(w)-\mathcal{E}_{\theta,X}^{\varphi}(w').
\end{equation}
In particular,
\begin{equation}\label{eq:LdiffonXsmoothmeasure2}
\lim_{k\to\infty}\mathcal{L}_{k,\varphi}(N^k_{w,\omega^n})=\mathcal{E}_{\theta,X}^{\varphi}(w)\,.
\end{equation}
\end{proposition}


\begin{proof}
First observe that by \cref{prop:BMNimplynorm}, for any $k\geq 1$, $N^k_{w,\omega^n}$ and $N^k_{w',\omega^n}$ are both norms, hence the expressions inside the limit in \eqref{eq:LdiffonXsmoothmeasure} make sense.

We first assume that $w,w'\in C^\infty(X)$. To start, we make the following observation:
\[
\begin{aligned}
\mathcal{L}_{k,\varphi}(N^k_{w,\omega^n}) - \mathcal{L}_{k,\varphi}(N^k_{w',\omega^n})&= \int_0^1 \frac{\mathrm{d}}{\mathrm{d}t} \mathcal{L}_{k,\varphi}(N^k_{w' + t(w-w'),\omega^n})\,\mathrm{d}t\\
& = \int_0^1 \int_X (w-w') \,\beta^k_{w' + t(w-w'),\varphi,\omega^n}\, \mathrm{d}t.
\end{aligned}
\]
The second equality follows from differentiating the Gram matrix of the Hilbert norm $N^k_{w'+t(w-w'),\omega^n}$:
\[
    \frac{\mathrm d}{\mathrm dt}\mathcal L_{k,\varphi}
    \left(N^k_{w'+t(w-w'),\omega^n}\right)=\int_X (w-w')\,\beta^k_{w'+t(w-w'),\varphi,\omega^n}.
\]

By \cref{prop:smooth_weak_conv}, we have
\[
\lim_{k\to\infty}\int_X (w-w') \,\beta^k_{w' + t(w-w'),\varphi,\omega^n}=\int_X (w-w')\,\theta_{P_{\theta,X}[\varphi](w' + t(w-w'))}^n.
\]
By \cref{thm:DXmain1}, we have
\[
\bigl|\int_X (w-w')\beta^k_{w' + t(w-w'),\varphi,\omega^n}\bigr| \leq C \sup_X |w - w'|.
\]
Hence, by the dominated convergence theorem we obtain that
\[
\begin{aligned}
\lim_{k\to\infty} \left(\mathcal{L}_{k,\varphi}(N^k_{w,\omega^n}) - \mathcal{L}_{k,\varphi}(N^k_{w',\omega^n})\right) &= \int_0^1 \int_X (w-w') \,\theta_{P_{\theta,X}[\varphi](w' + t(w-w'))}^n \,\mathrm{d}t\\
& = \mathcal{E}_{\theta,X}^{\varphi}(w)-\mathcal{E}_{\theta,X}^{\varphi}(w'),
\end{aligned}
\]
where in the last line we have used \cref{prop:differential_P}.

This proves \eqref{eq:LdiffonXsmoothmeasure} for smooth weights.

For general $w,w'\in C^0(X)$, choose smooth functions $w_j,w'_j$ converging uniformly to $w,w'$, respectively. If $\|w_j-w\|_{C^0}\leq \delta_j$, then
\[
\mathrm{e}^{-k\delta_j/2}N^k_{w,\omega^n}\leq N^k_{w_j,\omega^n}\leq \mathrm{e}^{k\delta_j/2}N^k_{w,\omega^n},
\]
and hence
\[
\left|\mathcal{L}_{k,\varphi}(N^k_{w_j,\omega^n})-\mathcal{L}_{k,\varphi}(N^k_{w,\omega^n})\right|
\leq \frac{n!}{k^n}h^0\left(X,L^k\otimes T\otimes \mathcal{I}(k\varphi)\right)\delta_j.
\]
The same estimate holds for $w'_j$ and $w'$. Moreover, \cref{prop:differential_P} shows that the energy side is Lipschitz with respect to the uniform norm of the weight. Letting first $k\to\infty$ in the smooth case, and then $j\to\infty$, proves \eqref{eq:LdiffonXsmoothmeasure} for continuous weights. Equation \eqref{eq:LdiffonXsmoothmeasure2} again follows by taking $w'=0$.
\end{proof}

\begin{lemma}\label{lma:BML}
Let $\varphi\in \PSH(X,\theta)$.
Let $(K,v)$ be a weighted subset of $X$. Let $\mu\in \BM(K,v)$. Then
\begin{equation}\label{eq:Bern_Mark_implies}
\lim_{k\to\infty} \big( \mathcal{L}_{k,\varphi}(N^k_{v,K}) - \mathcal{L}_{k,\varphi}(N^k_{v,\mu})\big) = 0\,.
\end{equation}
\end{lemma}
\begin{proof}
By \eqref{eq:Nkinfcomp} and \eqref{eq:BM}, for every $\epsilon>0$ and all $k\geq 1$,
\[
N^k_{v,\mu}\leq N^k_{v,K}\leq C_\epsilon \mathrm{e}^{\epsilon k}N^k_{v,\mu}.
\]
The corresponding unit balls therefore satisfy inclusions with radii differing by at most the factor $C_\epsilon \mathrm{e}^{\epsilon k}$. Taking logarithmic volumes and multiplying by $n!k^{-n-1}$ gives
\[
\left|\mathcal{L}_{k,\varphi}(N^k_{v,K})-\mathcal{L}_{k,\varphi}(N^k_{v,\mu})\right|
\leq \frac{2n!}{k^{n+1}} h^0\left(X,L^k\otimes T\otimes \mathcal{I}(k\varphi)\right)(\log C_\epsilon+\epsilon k).
\]
The right-hand side has limsup at most $2\epsilon\vol(\theta_\varphi)$ by \cref{thm:DXmain1}. Letting $\epsilon\to 0+$ proves \eqref{eq:Bern_Mark_implies}.
\end{proof}



\begin{corollary}\label{cor:Ninfdifflim}
Let $w \in C^0(X)$, $\varphi \in \PSH(X,\theta)$ be a potential with analytic singularities such that $\theta_\varphi$ is a Kähler current. Then
\[
\lim_{k\to\infty}\mathcal{L}_{k,\varphi}(N^k_{w,X})=\mathcal{E}_{\theta,X}^{\varphi}(w).
\]
\end{corollary}
\begin{proof}
This follows from \cref{lma:BML} and \cref{prop:quant_I_smooth}: indeed $\omega^n \in \BM(X,0)$, hence $\omega^n\in \BM(X,w)$ as $w$ is bounded.
\end{proof}


It would be helpful to consider the following auxiliary functions:
\[
\begin{aligned}
    P'_{\theta,K}[\varphi](v)\coloneqq & \sup \Bigl\{ \eta \in \PSH(X,\theta): \eta|_K \leq v  \textup{ and } \eta\preceq \varphi\Bigr\},\\
    P'_{\theta,K}[\varphi]_{\mathcal{I}}(v)\coloneqq & \sup \Bigl\{ \eta \in \PSH(X,\theta): \eta|_K \leq v  \textup{ and } \eta\preceq_{\mathcal{I}} \varphi\Bigr\}.
\end{aligned}
\]
\index{$P'_{\theta,K}[\varphi](v)$}\index{$P'_{\theta,K}[\varphi]_{\mathcal{I}}(v)$}

We note the following maximum principles that follow from the above definitions:
\begin{lemma}\label{lma:max_princ} Let $v\in C^0(K)$, $\varphi\in \PSH(X,\theta)$, and $\eta \in \PSH(X,\theta)$.
Assume that $\eta \preceq \varphi$. Then
\begin{equation}\label{eq:max_princ}
\sup_K (\eta-v) = \sup_{\{\eta\neq -\infty\}} \left(\eta - P'_{\theta,K}[\varphi](v)\right)=\sup_{\left\{P'_{\theta,K}[\varphi](v)\neq -\infty\right\}} \left(\eta - P'_{\theta,K}[\varphi](v)\right).
\end{equation}
\end{lemma}
\begin{proof}
We prove the first equality first. We write $S=\{\eta=-\infty\}$.

By definition, $P'_{\theta,K}[\varphi](v)|_K\leq v$, so
\[
\left.\left(\eta-P'_{\theta,K}[\varphi](v)\right)\right|_{K\setminus S}\geq \eta|_{K\setminus S}-v|_{K\setminus S}\,.
\]
This implies that
\[
\sup_K (\eta-v) \leq \sup_{X\setminus S} \left(\eta - P'_{\theta,K}[\varphi](v)\right).
\]

Conversely, observe that $\sup_K(\eta-v)>-\infty$ as $K$ is non-pluripolar. It follows that $\eta-\sup_K (\eta-v)$ is a candidate in the definition of $P'_{\theta,K}[\varphi](v)$. Hence,
\begin{equation}\label{eq:etaminussupK_temp1}
\eta-\sup_K (\eta-v)\leq P'_{\theta,K}[\varphi](v).
\end{equation}
This implies that
\[
\sup_K (\eta-v) \geq \sup_{X\setminus S} \left(\eta - P'_{\theta,K}[\varphi](v)\right),
\]
finishing the proof of the first identity in \eqref{eq:max_princ}.

In order to prove the latter identity, observe that
\[
\left\{P'_{\theta,K}[\varphi](v) =-\infty\right\} \subseteq S
\]
as a consequence of \eqref{eq:etaminussupK_temp1}. On $\{P'_{\theta,K}[\varphi](v)\neq -\infty\}\cap S$, we have $\eta-P'_{\theta,K}[\varphi](v)=-\infty$. Hence the last equality of \eqref{eq:max_princ} also follows.
\end{proof}

\begin{proposition}\label{prop:quant_I_algebraic_BM} Let $\varphi\in \PSH(X,\theta)$ be a potential with analytic singularities such that $\theta_{\varphi}$ is a Kähler current. Let $(K,v)$, $(K',v')$ be two weighted subsets of $X$.
Then
\begin{equation}\label{eq:LkdiffconvtoI}
\lim_{k\to\infty} \left( \mathcal{L}_{k,\varphi}(N^k_{v,K}) - \mathcal{L}_{k,\varphi}(N^k_{v',K'}) \right) = \mathcal{E}_{\theta,K}^{\varphi}(v)-\mathcal{E}_{\theta,K'}^{\varphi}(v').
\end{equation}
In particular,
\begin{equation}\label{eq:Lkconv}
\lim_{k\to\infty}\mathcal{L}_{k,\varphi}(N^k_{v,K})=\mathcal{E}_{\theta,K}^{\varphi}(v).
\end{equation}
\end{proposition}

\begin{proof}
First observe that by \cref{prop:BMNimplynorm}, for any $k>0$, $N^k_{v,K}$ and $N^k_{v',K'}$ are both norms, hence the expressions inside the limit in \eqref{eq:LkdiffconvtoI} make sense. Moreover, \eqref{eq:Lkconv} is just a special case of \eqref{eq:LkdiffconvtoI} for $K'=X$ and $v'= 0$.

To prove \eqref{eq:LkdiffconvtoI} it is enough to show that for any fixed $w \in C^\infty(X)$ we have
\begin{equation}\label{eq:inproofLdiffsmw}
\lim_{k\to\infty} \big( \mathcal{L}_{k,\varphi}(N^k_{v,K}) - \mathcal{L}_{k,\varphi}(N^k_{w,X}) \big) = \mathcal{E}_{\theta,K}^{\varphi}(v) - \mathcal{E}_{\theta,X}^{\varphi}(w)\,.
\end{equation}


For $\epsilon \in (0,1)$ small enough, $\theta_{(1-\epsilon)\varphi}$ is still a Kähler current. Let us fix such $\epsilon$, along with an arbitrary $\epsilon' \in (0,1)$.

Let $(v^-_{j})_j$ and $(v^+_{j})_j$ be the sequences of smooth functions constructed in \cref{lma:global_env_approx} for the data $(K,v)$.

By \cref{prop:Bergman_approx} there exists $k_0(\epsilon,\epsilon') \in \mathbb N$ such that
\[
\frac{1}{k} \log h^k \otimes h_T (s,s) \preceq (1-\epsilon)\varphi,
\]
and $\frac{1}{k} \log h^k \otimes h_T (s,s) \in \PSH(X,\theta + \epsilon'\omega)$ for any $s \in \mathrm{H}^0(X,T\otimes L^k \otimes \mathcal{I}(k\varphi))$, as long as $k \geq k_0(\epsilon,\epsilon')$.

In particular, \cref{lma:max_princ} gives that
\[
\begin{aligned}
N^k_{P'_{\theta + \epsilon'\omega,K}[(1-\epsilon)\varphi](v),X}(s)=& N^k_{v,K}(s),\\
N^k_{P'_{\theta + \epsilon'\omega,X}[(1-\epsilon)\varphi](v_{j}^-),X}(s)=& N^k_{v_{j}^-,X}(s),\\
N^k_{P'_{\theta + \epsilon'\omega,X}[(1-\epsilon)\varphi](v_{j}^+),X}(s)=& N^k_{v^+_{j},X}(s).
\end{aligned}
\]
As
\[
P'_{\theta + \epsilon'\omega,X}[(1-\epsilon)\varphi](v_{j}^-) \leq P'_{\theta + \epsilon'\omega,K}[(1-\epsilon)\varphi](v) \leq P'_{\theta + \epsilon'\omega,X}[(1-\epsilon)\varphi](v_{j}^+),
\]
where the second inequality follows from \cref{lma:global_env_approx} and the fact that $P'=P$ for the global continuous obstacles $v_j^\pm$. By \cref{lma:mononorm} we have
\[
N^k_{v^+_{j},X}(s) \leq N^k_{v,K}(s) \leq N^k_{v_{j}^-,X}(s), \quad s \in \mathrm{H}^0(X,T \otimes L^k  \otimes \mathcal{I}(k\varphi)), k \geq k_0(\epsilon,\epsilon').
\]
Composing with $\mathcal{L}_{k,\varphi}$ we arrive at
\[
\mathcal{L}_{k,\varphi}(N^k_{v^-_{j},X})\leq  \mathcal{L}_{k,\varphi}(N^k_{v,K})\leq  \mathcal{L}_{k,\varphi}(N^k_{v^+_{j},X})\,, \quad k \geq k_0(\epsilon,\epsilon').
\]
For any $j>0$, by \cref{cor:Ninfdifflim} we get
\[
\begin{aligned}
\mathcal{E}_{\theta,X}^{\varphi}(v^-_{j}) - \mathcal{E}_{\theta,X}^{\varphi}(w)=& \lim_{k\to\infty} \left(\mathcal{L}_{k,\varphi}(N^k_{v^-_{j},X})-\mathcal{L}_{k,\varphi}(N^k_{w,X})\right)\\
\leq & \varliminf_{k\to\infty}\left(\mathcal{L}_{k,\varphi}(N^k_{v,K})- \mathcal{L}_{k,\varphi}(N^k_{w,X})\right)\\
\leq & \varlimsup_{k\to\infty}\left(\mathcal{L}_{k,\varphi}(N^k_{v,K})- \mathcal{L}_{k,\varphi}(N^k_{w,X})\right)\\
\leq & \lim_{k\to\infty} \left(\mathcal{L}_{k,\varphi}(N^k_{v^+_{j},X})-\mathcal{L}_{k,\varphi}(N^k_{w,X})\right)\\
=&  \mathcal{E}_{\theta,X}^{\varphi}(v^+_{j}) - \mathcal{E}_{\theta,X}^{\varphi}(w).
\end{aligned}
\]
Using \cref{lma:global_env_approx}, we can let $j \to \infty$ to arrive at
\[
\begin{split}
\mathcal{E}_{\theta,K}^{\varphi}(v)  -\mathcal{E}_{\theta,X}^{\varphi}(w) &\leq \varliminf_{k\to\infty}\left(\mathcal{L}_{k,\varphi}(N^k_{v,K})- \mathcal{L}_{k,\varphi}(N^k_{w,X})\right)\\
&\leq\varlimsup_{k\to\infty}\left(\mathcal{L}_{k,\varphi}(N^k_{v,K})- \mathcal{L}_{k,\varphi}(N^k_{w,X})\right)\\
& \leq \mathcal{E}_{\theta,K}^{\varphi}(v)  -\mathcal{E}_{\theta,X}^{\varphi}(w).
\end{split}
\]
Hence, \eqref{eq:inproofLdiffsmw} follows.
\end{proof}






\begin{corollary} \label{cor:LktoI}
Let $\varphi \in \PSH(X,\theta)$ be a potential with analytic singularities such that $\theta_{\varphi}$ is a Kähler current.
Then
\[
\lim_{k\to\infty}  \mathcal{L}_{k,\varphi}(N^k_{v,\mu})= \mathcal{E}_{\theta,K}^{\varphi}(v).
\]
\end{corollary}
\begin{proof}
Our claim follows from \cref{prop:quant_I_algebraic_BM} and \cref{lma:BML}.
\end{proof}


\begin{proposition}\label{prop:weakconvana}
Let $\varphi \in \PSH(X,\theta)$ be a potential with analytic singularities such that $\theta_{\varphi}$ is a Kähler current. Then
\[
\beta^k_{v,\varphi, \mu} \rightharpoonup \theta_{P_{\theta,K}[\varphi]_{\mathcal{I}}(v)}^n=\theta_{P_{\theta,K}[\varphi](v)}^n\quad \textup{as }k\to\infty.
\]
\end{proposition}

\begin{proof}
For $w\in C^0(X)$, let $f_k,g\colon \mathbb{R}\rightarrow \mathbb{R}$ ($k\geq 1$) be defined as
\[
f_k(t)\coloneqq \mathcal{L}_{k,\varphi}(N^k_{v+tw,\mu}),\quad g(t)\coloneqq \mathcal{E}_{\theta,K}^{\varphi}(v+tw).
\]
By \cref{cor:LktoI}, for all $t\in \mathbb{R}$, we have
\[
\lim_{k\to\infty} f_k(t)= g(t).
\]
Here we use that $\mu\in\BM(K,v+tw)$ for all fixed $t$, which follows from $\mu\in\BM(K,v)$ and the boundedness of $tw$ on $K$. Note that $f_k$ is concave by H\"older's inequality.
For each $k$, the function $f_k$ is differentiable and
\[
    f_k'(0)=\int_X w\,\beta^k_{v,\varphi,\mu}.
\]
This follows by differentiating the Gram matrix of the Hilbert norm $N^k_{v+tw,\mu}$. The function $g$ is differentiable at $0$ by \cref{prop:differential_P}, with
\[
    g'(0)=\int_K w\,\theta^n_{P_{\theta,K}[\varphi]_{\mathcal I}(v)}.
\]
Applying \cref{lma:conc_lim_der} to the convex functions $-f_k$ and $-g$ gives
\[
    \lim_{k\to\infty} f_k'(0)=g'(0),
\]
which is equivalent to $\beta^k_{v, \varphi,\mu} \rightharpoonup \theta_{P_{\theta,K}[\varphi](v)}^n$, by \cref{prop:differential_P}.
\end{proof}


\begin{proposition} \label{prop:mainKahcurr}
Suppose that $\varphi \in \PSH(X,\theta)$ is such that $\theta_{\varphi}$ is a Kähler current.
Let $(K,v)$ be a weighted subset of $X$ and $\mu\in \BM(K,v)$. Then
\begin{equation}\label{eq:pbkconv1}
\beta^k_{v,\varphi,\mu} \rightharpoonup \theta_{P_{\theta,K}[\varphi]_{\mathcal{I}}(v)}^n\quad \textup{as } k\to\infty.
\end{equation}
\end{proposition}
\begin{proof} Let $\tau$ be the weak limit of a subsequence of $\beta^k_{v,\varphi,\mu}$. We claim that
\begin{equation}\label{eq:inproofmuleq}
\tau\leq \theta_{P_{\theta,K}[\varphi]_{\mathcal{I}}(v)}^n.
\end{equation}
Observe that this claim implies the conclusion. In fact, by \cref{thm:DXmain1}, we have equality of the total masses, so equality holds in \eqref{eq:inproofmuleq}. As $\tau$ is an arbitrary cluster point of the sequence $(\beta^k_{v,\varphi,\mu})_k$, we get \eqref{eq:pbkconv1}.

It remains to prove \eqref{eq:inproofmuleq}.
Let $(\varphi_j)$ be a quasi-equisingular approximation of $\varphi$ in $\PSH(X,\theta)$. We may assume that $\theta_{\varphi_j}$ is a K\"ahler current for all $j\geq 1$.
By \cref{cor:quasi-equichar}, we know that
\[
\varphi_j\xrightarrow{d_S}P_{\theta}[\varphi]_{\mathcal{I}}.
\]
The quasi-equisingular approximation is decreasing, so applying the decreasing case of \cref{prop:conv_of_K_env} to the limit $P_{\theta}[\varphi]_{\mathcal{I}}$, and using the definition of the relative $\mathcal{I}$-envelope, we get
\[
P_{\theta,K}\left[\varphi_j\right]_{\mathcal{I}}(v)\searrow P_{\theta,K}[\varphi]_{\mathcal{I}}(v).
\]
In particular,
\begin{equation}\label{eq:inproofeqmass}
\lim_{j\to\infty} \int_X \theta_{P_{\theta,K}\left[\varphi_j\right]_{\mathcal{I}}(v)}^n= \int_X  \theta_{P_{\theta,K}[\varphi]_{\mathcal{I}}(v)}^n.
\end{equation}
Observe that
\[
\beta^k_{v,\varphi,\mu} \leq \beta^k_{v,\varphi_j,\mu}
\]
for any $k \geq 1$.
As $\mu\in \BM(K,v)$,
by \cref{prop:weakconvana},
\[
\tau \leq \theta_{P_{\theta,K}[\varphi_j]_{\mathcal{I}}(v)}^n,
\]
for any $j \geq 1$ fixed.
By \eqref{eq:inproofeqmass} and the semicontinuity theorem \cref{thm:semicon}, we have
\[
\theta_{P_{\theta,K}\left[\varphi_j\right]_{\mathcal{I}}(v)}^n\rightharpoonup \theta_{P_{\theta,K}[\varphi]_{\mathcal{I}}(v)}^n
\]
as $j\to\infty$. Therefore, \eqref{eq:inproofmuleq} follows.
\end{proof}





\begin{proof}[Proof of \cref{thm:pBMconvergence}]
By \cref{lma:same_sing_type} and \cref{prop:Ienvprojection}, the three potentials below have the same multiplier ideals. Hence
\[
\begin{aligned}
\mathrm{H}^0\left(X,L^k \otimes T \otimes \mathcal{I}(k\varphi)\right) &= \mathrm{H}^0\left(X,L^k \otimes T \otimes \mathcal I(kP_{\theta}[\varphi]_{\mathcal{I}})\right)\\
 &= \mathrm{H}^0\left(X,L^k \otimes T \otimes \mathcal I\left(kP_{\theta,K}[\varphi]_{\mathcal{I}}(v)\right)\right).
 \end{aligned}
\]
This allows us to replace $\varphi$ with $P_{\theta,K}[\varphi]_{\mathcal{I}}(v)$.


By \cref{lma:kahcurrentposmass}, there exists a sequence $(\varphi_j)_j$ in $\PSH(X,\theta)$ such that $\varphi_j\nearrow \varphi$ a.e.\ and $\theta_{\varphi_j}$ is a Kähler current for each $j\geq 1$. This gives
\[
\beta^k_{v,\varphi_j,\mu} \leq \beta^k_{v,\varphi,\mu}.
\]
Let $\tau$ be the weak limit of a subsequence of $(\beta^k_{v,\varphi,\mu})_k$. Then by \cref{prop:mainKahcurr},
\[
\theta_{P_{\theta,K} [\varphi_j]_{\mathcal{I}}(v)}^n\leq \tau.
\]
By \cref{prop:conv_of_K_env} and \cref{cor:incseqnppcont}, we have
\[
\theta_{P_{\theta,K} [\varphi_j]_{\mathcal{I}}(v)}^n \rightharpoonup \theta_{P_{\theta,K} [\varphi]_{\mathcal{I}}(v)}^n.
\]
Hence,
\begin{equation}\label{eq:inproofmulower}
\theta_{P_{\theta,K} [\varphi]_{\mathcal{I}}(v)}^n\leq \tau.
\end{equation}
A comparison of total masses using \eqref{eq:betaint_conv} gives that equality holds in \eqref{eq:inproofmulower}. As $\tau$ is arbitrary, we obtain \eqref{eq:pbkconvgeneral}.
\end{proof}

\begin{remark}\label{rmk:Bergman:L1290}
    The results in this chapter can also be reformulated as the large deviation principle of a determinantal point process on $X$ using the Gärtner--Ellis theorem exactly as in \cite{Ber14}. We do not pursue this reformulation here.
\end{remark}
